\chapter{Elliptic Curves}
\begin{center}\itshape
The right triangle with sides $3,4,5$ has area $6$.\\
So does the one with sides $\frac7{10},\frac{120}7,\frac{1201}{70}$.\\
Where did that come from, and is there a third?
\end{center}

Around 972, an anonymous Arab manuscript asked which whole numbers are the area of a right triangle with rational sides. Such numbers
are called \emph{congruent numbers}. In 1225 Leonardo of Pisa, Fibonacci, was challenged at the court of Frederick~II to show that $5$ is
one, and found the triangle $\frac32,\frac{20}3,\frac{41}6$. Four centuries later Fermat proved that $1$ is not, by his method of infinite
descent (Volume~1, Theorem~9.2): no right triangle with rational sides has a square area. Which numbers are congruent is still not fully
known; the answer, if a conjecture worth a million dollars is true, was found by Tunnell in 1983.

The problem is about a curve. A right triangle of area $n$ is the same thing as a rational point with $y\neq0$ on the cubic
\[
y^2=x^3-n^2x,
\]
and the triangle $\frac7{10},\frac{120}7,\frac{1201}{70}$ comes from the point $(12,36)$ that the $3$-$4$-$5$ triangle gives on $y^2=x^3-36x$, by
the oldest trick in the subject: draw the tangent at a rational point, and its third intersection with the cubic is rational again. Diophantus did this
around 250, Bachet and Fermat in the seventeenth century, Newton saw that chords and tangents make the rational points of a cubic into
something like an arithmetic, and Poincar\'e asked in 1901 whether finitely many points always generate all the others. Mordell proved it in 1922.

This chapter builds the arithmetic: the group law, proved associative with the nine-point theorem of Chapter~11; the torsion points, which are
integral (Nagell--Lutz) and inject into the curve modulo $p$; Mordell's theorem, by a descent that needs the finiteness of class groups and the unit
theorem of Chapters~3 and~4; and the congruent number problem, with a complete proof that $1$ is not congruent, hence Fermat's theorem
$x^4+y^4\ne z^4$ once more, now as the statement that a certain elliptic curve has rank $0$.

\begin{reminder}[from Chapter 11 and Volume 1]
A smooth plane cubic $C$ has genus $1$; for every point $O$, each divisor of degree $0$ is linearly equivalent to $P-O$ for a unique $P$, and two cubics
whose divisors on $C$ share eight points share the ninth (Corollary~\ref{cor:genus1}). A line meets a plane cubic not containing it in three points
counted with multiplicity (Corollary~\ref{cor:bezoutcons}), and Lemma~\ref{lem:finite} holds over any field (its proof uses only that $K[y]$ is a
PID and Gauss's lemma). \emph{Volume~1.} Fermat's theorem: $x^4+y^4=z^2$ has no solution in positive integers (Volume~1, Theorem~9.2); the right triangles
with integer sides (Volume~1, Theorem~9.1). In this chapter $K$ is a field of characteristic $\ne2,3$ and $\overline K$ an algebraic closure.
\end{reminder}

\section{The group law}

An \emph{elliptic curve} over $K$ is a curve
\[
E:\ y^2=x^3+ax+b,\qquad a,b\in K,\qquad 4a^3+27b^2\ne0,
\]
with its point at infinity $O=(0:1:0)$; projectively $Y^2Z=X^3+aXZ^2+bZ^3$. The condition says that $f(x)=x^3+ax+b$ has three distinct roots in
$\overline K$ (its discriminant is $-(4a^3+27b^2)$), and this is exactly smoothness: an affine singular point has $2y=0$ and $f'(x)=0$, so $x$ is a double
root of $f$; at infinity, $F_Z=Y^2+\dots=1\ne0$. By Exercise~\ref{ex:weierstrass}, the line at infinity $Z=0$ meets $E$ only at $O$, with multiplicity
$3$: $O$ is a \emph{flex}. Every smooth cubic with a rational flex can be brought to this \emph{Weierstrass form} by a change of coordinates over $K$
(send the flex to $(0:1:0)$ and its tangent to $Z=0$, then complete the square and the cube); we take it as the definition. For a field $L\supseteq K$,
$E(L)$ is the set of points with coordinates in $L$, together with $O$.

\subsection*{Chords and tangents}

For $P,Q\in E(\overline K)$, the line through $P$ and $Q$ (the tangent if $P=Q$) meets $E$ in a third point, counted with multiplicity
(Corollary~\ref{cor:bezoutcons}); call it $P*Q$. If $P,Q\in E(L)$ then $P*Q\in E(L)$: when $P,Q$ are affine and the line is $y=\lambda x+\nu$ with
$\lambda,\nu\in L$, substituting gives
\[
x^3+ax+b-(\lambda x+\nu)^2=(x-x_P)(x-x_Q)(x-x_{P*Q}),
\]
and comparing the coefficients of $x^2$, $x_{P*Q}=\lambda^2-x_P-x_Q\in L$. (The roots of this cubic are the intersection points with their
multiplicities, because along the line, $x$ is a uniformiser at each affine point; Proposition~\ref{prop:smoothdvr}.) Vertical lines $x=c$ pass through $O$
(the line $X=cZ$ contains $(0:1:0)$), so if $x_P=x_Q$ and $Q\ne P$, then $P*Q=O$.

\begin{definition}
$P+Q=O*(P*Q)$.
\end{definition}

Concretely, the line through $O$ and $R=(x,y)$ is the vertical $x=x_R$, so $O*R=(x,-y)$: to add $P$ and $Q$, draw the line, take the third point, and
reflect it in the $x$-axis (Figure~\ref{fig:chord}). In formulas, for $P=(x_1,y_1)$, $Q=(x_2,y_2)$ affine with $Q\ne-P$,
\begin{equation}\label{eq:addformula}
\begin{gathered}
\lambda=\begin{cases}\dfrac{y_2-y_1}{x_2-x_1}&(x_1\ne x_2)\\[2ex]\dfrac{3x_1^2+a}{2y_1}&(P=Q)\end{cases}\qq
x_3=\lambda^2-x_1-x_2,\quad y_3=-y_1-\lambda(x_3-x_1),
\end{gathered}
\end{equation}
and $P+Q=(x_3,y_3)$; also $P+O=P$ and $(x,y)+(x,-y)=O$.

\begin{figure}[ht]
\centering\fig{ch12-chord}
\caption{The group law on $y^2=x^3-x+1$. The line through $P=(-1,1)$ and $Q=(1,1)$ meets the curve again at $P*Q=(0,1)$; reflecting gives
$P+Q=(0,-1)$. The tangent at $P$ gives $2P=(3,-5)$.}
\label{fig:chord}
\end{figure}

\begin{think}
Why does the definition need $O$ at all? Try to define $P+Q$ as $P*Q$: is there a neutral element? And where is $O$ on Figure~\ref{fig:chord}: what
is the third point of a vertical line?
\end{think}

\begin{theorem}\label{thm:grouplaw}
$(E(L),+)$ is an abelian group with neutral element $O$, and $-(x,y)=(x,-y)$. Three points $P,Q,R$ satisfy $P+Q+R=O$ iff they are the intersection of
$E$ with a line (with multiplicities).
\end{theorem}
\begin{proof}
$E(L)$ is stable by the remarks above. Commutativity is clear (the line through $P,Q$ is the line through $Q,P$). $P+O=O*(P*O)=P$: the line through $P$ and
$O$ meets $E$ in $P,O,P*O$, so the third point of the line through $O$ and $P*O$ is $P$. The inverse of $P$ is $P'=O*P$: the line through $P$ and $P'$
passes through $O$ (since $O,P,P'$ are collinear), so $P*P'=O$, and $P+P'=O*O=O$, because the tangent at the flex $O$ meets $E$ only at $O$. If
$P,Q,R$ lie on a line, then $P*Q=R$ and $P+Q=O*R=-R$; conversely, if $P+Q=-R$, then $O*(P*Q)=O*R$, so $P*Q=R$.

\emph{Associativity.} Let $P,Q,R\in E(L)$. Introduce the points
\[
S'=P*Q,\ S=P+Q,\ T'=S*R,\qquad U'=Q*R,\ U=Q+R,\ T''=P*U,
\]
so that $(P+Q)+R=O*T'$ and $P+(Q+R)=O*T''$; it suffices to show $T'=T''$. Consider the six lines
\[
\ell_1=PQS',\quad\ell_2=SRT',\quad\ell_3=U'OU,\qquad m_1=S'OS,\quad m_2=QRU',\quad m_3=PUT'',
\]
each named by its three intersection points with $E$. The cubics $A=\ell_1\ell_2\ell_3$ and $B=m_1m_2m_3$ have
\[
\begin{gathered}
\operatorname{div}(A)=P+Q+S'+S+R+T'+U'+O+U,\\
\operatorname{div}(B)=S'+O+S+Q+R+U'+P+U+T''.
\end{gathered}
\]
They share the eight points $P,Q,R,S,S',U,U',O$ (as divisors), so by the nine-point theorem (Corollary~\ref{cor:genus1}(d)), $T'=T''$.
\end{proof}

The proof is Figure-free and multiplicity-safe because the nine-point theorem was proved for divisors. In particular it covers $P=Q$ and all the
degenerate configurations, which a proof by explicit formulas would have to check one by one.

\begin{corollary}\label{cor:pic}
The map $P\mapsto[P-O]$ is an isomorphism from $E(\overline K)$ to the group of divisor classes of degree $0$ on $E$. A divisor $\sum n_iP_i$ of degree $0$ is
principal iff $\sum n_iP_i=O$ in the group $E$.
\end{corollary}
\begin{proof}
If $\ell$ is the line through $P,Q,P*Q$ and $\ell'$ the line through $P*Q,O,P+Q$, then
$\operatorname{div}(\ell/\ell')=P+Q-O-(P+Q)$, so $(P+Q)-O\equiv(P-O)+(Q-O)$: the map is a homomorphism. It is bijective by Corollary~\ref{cor:genus1}(c).
For the second statement, $\sum n_iP_i=\sum n_i(P_i-O)\equiv(\sum n_iP_i)_E-O$, which is principal iff $(\sum n_iP_i)_E=O$ (Corollary~\ref{cor:genus1}(b)).
\end{proof}

\section{Torsion points}

From now on, $E:y^2=x^3+ax+b$ with $a,b\in\Z$, and $D=4a^3+27b^2\ne0$. A point $P$ has \emph{finite order} if $mP=O$ for some $m\ge1$; these form
the \emph{torsion subgroup} $E(\Q)_{\mathrm{tors}}$. For instance on $y^2=x^3+1$, $(-1,0)$ has order $2$, $(0,\pm1)$ have order $3$ (they are flexes,
Exercise~\ref{ex:flex}) and $(2,\pm3)$ have order $6$.

\subsection*{Denominators}

Let $p$ be a prime and $P=(x,y)\in E(\Q)$ with $v_p(x)<0$. Then $v_p(x^3+ax+b)=3v_p(x)$, so $2v_p(y)=3v_p(x)$, and there is an integer $\nu\ge1$ with
$v_p(x)=-2\nu$, $v_p(y)=-3\nu$. Write $x=m/e^2$, $y=n/e^3$ with $e\ge1$ and $\gcd(m,e)=\gcd(n,e)=1$: the denominators of rational points are squares and
cubes. For $\nu\ge1$ let
\[
E_\nu=\{P\in E(\Q):v_p(x(P))\le-2\nu\}\cup\{O\},
\]
the points that are $p$-adically close to $O$. Near $O$ use the coordinates $t=x/y$, $s=1/y$ (the chart $Y=1$), in which $O=(0,0)$ and
\[
E:\ s=t^3+ats^2+bs^3.
\]
For $P\ne O$ in $E_\nu$, with $v_p(x)=-2\nu'$ ($\nu'\ge\nu$), we have $v_p(t)=\nu'$ and $v_p(s)=3\nu'$; so $P\in E_{\nu''}$ iff $v_p(t(P))\ge\nu''$. Conversely,
any point with $v_p(s)\ge3\nu$ lies in $E_\nu$: then $v_p(y)<0$, which forces $v_p(x)<0$ (if $v_p(x)\ge0$ then $y^2=f(x)$ is $p$-integral), and
$v_p(y)=-3\nu'$ with $3\nu'\ge3\nu$. We put $t(O)=s(O)=0$.

\begin{proposition}[the $p$-adic filtration]\label{prop:filtration}
Each $E_\nu$ is a subgroup of $E(\Q)$, and $t$ induces an injective homomorphism $E_\nu/E_{3\nu}\to p^\nu\Z_{(p)}/p^{3\nu}\Z_{(p)}$. Consequently $E_1$
contains no point of finite order other than $O$.
\end{proposition}
\begin{proof}
\emph{The curve is a graph over $t$.} Let $P_1,P_2\in E_\nu$ have the same $t$. Subtracting their equations,
$s_2-s_1=(s_2-s_1)\bigl[at(s_1+s_2)+b(s_1^2+s_1s_2+s_2^2)\bigr]$, and the bracket has positive valuation, so $s_1=s_2$: $P_1=P_2$.

\emph{The line.} Let $P_1,P_2\in E_\nu$, not both $O$, and let $s=\alpha t+\beta$ be the line through them in the $(t,s)$-plane (the tangent if $P_1=P_2$; by the
previous point it is not vertical). If $t_1\ne t_2$, subtract the two equations and divide by $t_2-t_1$, using
$t_2s_2^2-t_1s_1^2=(t_2-t_1)s_2^2+t_1(s_2-s_1)(s_2+s_1)$:
\[
\alpha=\frac{s_2-s_1}{t_2-t_1}=\frac{t_1^2+t_1t_2+t_2^2+as_2^2}{1-at_1(s_1+s_2)-b(s_1^2+s_1s_2+s_2^2)}.
\]
The same formula with $t_1=t_2$, $s_1=s_2$ gives the slope of the tangent (implicit differentiation). The denominator is a $p$-adic unit, so
$v_p(\alpha)\ge2\nu$, and $\beta=s_1-\alpha t_1$ has $v_p(\beta)\ge3\nu$.

\emph{The third point.} Substituting $s=\alpha t+\beta$ in the equation gives a cubic in $t$ whose roots are $t_1,t_2,t_3$, the $t$-coordinates of the three
intersection points:
\[
(1+a\alpha^2+b\alpha^3)\,t^3+(2a\alpha\beta+3b\alpha^2\beta)\,t^2+\dots=0,
\qquad t_1+t_2+t_3=-\frac{2a\alpha\beta+3b\alpha^2\beta}{1+a\alpha^2+b\alpha^3}.
\]
So $v_p(t_1+t_2+t_3)\ge5\nu$; in particular $v_p(t_3)\ge\nu$ and $s_3=\alpha t_3+\beta$ has $v_p(s_3)\ge3\nu$, so the third point $P_3$ is in $E_\nu$ (or is $O$). Since
$P_1+P_2=-P_3$ and $-(x,y)=(x,-y)$ changes the signs of $t$ and $s$,
\[
t(P_1+P_2)=-t_3\equiv t(P_1)+t(P_2)\pmod{p^{5\nu}}.
\]
Together with $t(-P)=-t(P)$, this shows that $E_\nu$ is a subgroup and that $t$ is a homomorphism modulo $p^{5\nu}$, hence modulo $p^{3\nu}$. Its
kernel modulo $p^{3\nu}$ is $\{P:v_p(t(P))\ge3\nu\}=E_{3\nu}$.

\emph{No torsion.} Let $P\in E_1$ have finite order $m>1$, and $P\in E_\nu\setminus E_{\nu+1}$, so $v_p(t(P))=\nu$. If $p\nmid m$, then
$0=t(mP)\equiv m\,t(P)\pmod{p^{5\nu}}$, which is absurd since $v_p(m\,t(P))=\nu<5\nu$. If $p\mid m$, replace $P$ by $P'=\frac mpP$, a point of order
$p$, in some $E_{\nu'}\setminus E_{\nu'+1}$ ($E_1$ is a group): $0=t(pP')\equiv p\,t(P')\pmod{p^{5\nu'}}$ with $v_p(p\,t(P'))=\nu'+1<5\nu'$. Absurd again.
\end{proof}

\begin{theorem}[Nagell--Lutz, 1935--1937]\label{thm:nagelllutz}
If $P=(x,y)\in E(\Q)$ has finite order, then $x,y\in\Z$, and either $y=0$ or $y^2\mid4a^3+27b^2$.
\end{theorem}
\begin{proof}
For every prime $p$, $P\notin E_1$ by Proposition~\ref{prop:filtration}, so $v_p(x)\ge0$: $x\in\Z$, and $y^2=f(x)\in\Z$ forces $y\in\Z$. Suppose $y\ne0$. Then
$2P\ne O$ is also of finite order, so $x(2P)\in\Z$. By \eqref{eq:addformula}, $x(2P)=\lambda^2-2x$ with $\lambda=f'(x)/2y$; so $\lambda^2\in\Z$, hence
$\lambda\in\Z$, and $f'(x)=2y\lambda$ is divisible by $y$. The polynomial identity
\[
4a^3+27b^2=(3x^2+4a)\,f'(x)^2-27(x^3+ax-b)\,f(x),
\]
which one checks by expanding, then shows $y^2\mid4a^3+27b^2$, since $y^2\mid f'(x)^2$ and $y^2=f(x)$.
\end{proof}

So the torsion is finite and can be listed: try the finitely many $y$ with $y=0$ or $y^2\mid D$, solve for integral $x$, and test which points have finite
order. Many curves are disposed of at once: on $y^2=x^3+2$, $D=108$, and the candidates $y\in\{\pm1,\pm2,\pm3,\pm6\}$ give no integral $x$ except
$y=\pm1$, $x=-1$; and $P=(-1,1)$ has $2P=(\frac{17}4,-\frac{71}8)$, not integral, so $P$ has infinite order. The torsion of $y^2=x^3+2$ is trivial, and its rank is
at least $1$.

\subsection*{Reduction modulo $p$}

For $p\nmid2D$, the reduced curve $\tilde E:y^2=x^3+\bar ax+\bar b$ over $\mathbb F_p$ is an elliptic curve. Reduction sends $P=(x,y)\in E(\Q)$ to
$\tilde P=(\bar x,\bar y)$ if $v_p(x)\ge0$, and to $\tilde O$ if $P\in E_1$ or $P=O$.

\begin{proposition}\label{prop:reduction}
For $p\nmid2D$, reduction $E(\Q)\to\tilde E(\mathbb F_p)$ is a homomorphism with kernel $E_1$. It is injective on $E(\Q)_{\mathrm{tors}}$; in particular
$|E(\Q)_{\mathrm{tors}}|$ divides $|\tilde E(\mathbb F_p)|$.
\end{proposition}
\begin{proof}
In projective coordinates, write each point of $E(\Q)$ as $(X:Y:Z)$ with coprime integers; reduction is then $(X:Y:Z)\mapsto(\bar X:\bar Y:\bar Z)$, which
agrees with the description above. Let $P,Q,R$ be the intersection of $E$ with a line $\ell:uX+vY+wZ=0$, $u,v,w$ coprime integers. Choose a basis
$A,B\in\Z^3$ of the lattice $\ell\cap\Z^3$ (it has rank $2$, and is saturated: $\Z^3/(\ell\cap\Z^3)$ is torsion-free), so that every point of $\ell$ is
$\sigma A+\tau B$, $(\sigma:\tau)\in\mathbb P^1$; the reductions $\bar A,\bar B$ are independent (they span the plane $\bar u X+\bar vY+\bar wZ=0$, since $A,B$
extend to a basis of $\Z^3$). The binary cubic $\varphi(\sigma,\tau)=F(\sigma A+\tau B)$ has integer coefficients and factors over $\Q$ as
$c\prod_{P}(\tau_P\sigma-\sigma_P\tau)$ over the three points with their multiplicities, with $(\sigma_P,\tau_P)$ coprime integers; by Gauss's lemma, $c\in\Z$.
Reduce: $\bar\varphi=\bar c\prod(\bar\tau_P\sigma-\bar\sigma_P\tau)$, and $\bar c\ne0$ because $\tilde E$ is irreducible and so does not contain the line
$\bar\ell$. The roots of $\bar\varphi$ are the intersection points of $\bar\ell$ with $\tilde E$, with multiplicities (as in the affine argument above), so they are
$\tilde P,\tilde Q,\tilde R$, and $\tilde P+\tilde Q+\tilde R=\tilde O$ by Theorem~\ref{thm:grouplaw}. Since $\tilde O$ reduces from $O$ and $-P$ reduces
to $-\tilde P$, this shows that reduction is a homomorphism. Its kernel is $E_1\cup\{O\}=E_1$, which has no non-trivial torsion
(Proposition~\ref{prop:filtration}).
\end{proof}

For example, for $E:y^2=x^3+1$, $D=27$, and $\tilde E(\mathbb F_5)$ and $\tilde E(\mathbb F_{11})$ have $6$ and $12$ points (Exercise~\ref{ex:supersing}); so
$|E(\Q)_{\mathrm{tors}}|$ divides $6$, and the six points listed above are all of it.

\begin{remark}
Mazur proved in 1977 that $E(\Q)_{\mathrm{tors}}$ is one of $\Z/m$ for $1\le m\le10$ or $m=12$, or $\Z/2\times\Z/2m$ for $1\le m\le4$, and all fifteen occur.
The proof uses modular curves, the objects of Chapter~14.
\end{remark}

\section{Mordell's theorem}

\begin{theorem}[Mordell, 1922]\label{thm:mordell}
$E(\Q)$ is a finitely generated abelian group: $E(\Q)\cong E(\Q)_{\mathrm{tors}}\times\Z^r$. The integer $r$ is the \emph{rank} of $E$.
\end{theorem}

The proof has two halves: $E(\Q)/2E(\Q)$ is finite (the \emph{weak} Mordell--Weil theorem), and a \emph{height} function measures the size of
points in a way that halving decreases. Fermat's infinite descent is the model: a point is written as $Q_i+2P_1$ with $Q_i$ from a finite list and
$P_1$ smaller, and so on until the points are too small to go on.

\subsection*{The weak Mordell--Weil theorem}

Write $f(x)=(x-e_1)(x-e_2)(x-e_3)$ with $e_i\in\overline{\Q}$, distinct algebraic integers, and let $K=\Q(e_1,e_2,e_3)$, a number field of degree $1,2,3$ or $6$,
over which the three points $T_i=(e_i,0)$ of order $2$ are rational. For $i=1,2,3$ and $\{i,j,k\}=\{1,2,3\}$ define
$\delta_i:E(K)\to K^\times/K^{\times2}$ by
\[
\delta_i(x,y)=x-e_i\ (x\ne e_i),\qquad\delta_i(T_i)=(e_i-e_j)(e_i-e_k),\qquad\delta_i(O)=1.
\]

\begin{lemma}\label{lem:deltahom}
Each $\delta_i$ is a homomorphism, and $P\in2E(K)$ iff $\delta_1(P),\delta_2(P),\delta_3(P)$ are all squares.
\end{lemma}
\begin{proof}
\emph{Homomorphism.} As $\delta_i(-P)=\delta_i(P)$ and the target has exponent $2$, it suffices to show that $\delta_i(P)\delta_i(Q)\delta_i(R)$ is a square
when $P+Q+R=O$. If one of them is $O$, the other two are opposite and the product is $\delta_i(P)^2$. Otherwise $P,Q,R$ lie on a non-vertical line
$y=\lambda x+\nu$, and $f(x)-(\lambda x+\nu)^2=(x-x_P)(x-x_Q)(x-x_R)$. At $x=e_i$:
\[
(x_P-e_i)(x_Q-e_i)(x_R-e_i)=(\lambda e_i+\nu)^2.
\]
If none of the three points is $T_i$, the right side is non-zero and we are done. If $P=T_i$, then $\lambda e_i+\nu=0$, so $\lambda x+\nu=\lambda(x-e_i)$ and
$(x-x_Q)(x-x_R)=(x-e_j)(x-e_k)-\lambda^2(x-e_i)$; at $x=e_i$, $(x_Q-e_i)(x_R-e_i)=(e_i-e_j)(e_i-e_k)=\delta_i(T_i)$. Neither $Q$ nor $R$ is $T_i$ (the line would
be tangent at $T_i$, whose tangent is vertical), so $\delta_i(Q)\delta_i(R)=\delta_i(T_i)$ and the product of the three is a square.

\emph{Kernel.} If $P=2Q$, then $\delta_i(P)=\delta_i(Q)^2$. Conversely, let $P=(x,y)$ with all $\delta_i(P)$ squares; then each $x-e_i$ is a square $r_i^2$ in $K$
(for $P=T_i$, $x-e_i=0$, and $x-e_j=e_i-e_j=\delta_j(T_i)$). We look for the tangent line $y=\lambda X+\nu$ at a point $Q$ with $2Q=\pm P$; it must satisfy
$f(X)-(\lambda X+\nu)^2=(X-x_Q)^2(X-x)$. Choose signs $\varepsilon_i=\pm1$ such that the linear system
\[
\varepsilon_ir_i\,(e_i-x_Q)=\lambda e_i+\nu\qquad(i=1,2,3)
\]
in the unknowns $x_Q,\lambda,\nu$ has a non-zero determinant $\sum_{\text{cyclic}}\varepsilon_ir_i(e_j-e_k)$: if it vanished for all signs, flipping
$\varepsilon_1$ would give $r_1(e_2-e_3)=0$, and similarly $r_2=r_3=0$, so $x=e_1=e_2=e_3$, absurd. Solve it in $K$. The polynomial
$h(X)=f(X)-(\lambda X+\nu)^2-(X-x_Q)^2(X-x)$ has degree $\le2$ and vanishes at the three $e_i$, since $(e_i-x_Q)^2(e_i-x)=-r_i^2(e_i-x_Q)^2=-(\lambda e_i+\nu)^2$.
So $h=0$, the line meets $E$ at $Q=(x_Q,\lambda x_Q+\nu)\in E(K)$ doubly and at a point with first coordinate $x$, which is $P$ or $-P$. Hence $2Q=-(\pm P)$, and
$P=2Q$ or $P=2(-Q)$.
\end{proof}

\begin{lemma}\label{lem:Sfinite}
Let $S$ be the set of primes of $\mathcal O_K$ dividing $2(e_1-e_2)(e_1-e_3)(e_2-e_3)$. Each $\delta_i$ takes values in the group
\[
K(S,2)=\{\alpha\in K^\times/K^{\times2}:\ v_{\mathfrak p}(\alpha)\text{ is even for every prime }\mathfrak p\notin S\},
\]
which is finite.
\end{lemma}
\begin{proof}
Let $P=(x,y)\in E(K)$, $P\ne T_i$, and $\mathfrak p\notin S$. If $v_{\mathfrak p}(x)<0$, then $v_{\mathfrak p}(x-e_i)=v_{\mathfrak p}(x)$ for all $i$ ($e_i$ is
integral), and $3v_{\mathfrak p}(x)=2v_{\mathfrak p}(y)$, so it is even. If $v_{\mathfrak p}(x)\ge0$, at most one of the $x-e_i$ is divisible by $\mathfrak p$
(otherwise $\mathfrak p$ divides some $e_i-e_j$), and $\sum_iv_{\mathfrak p}(x-e_i)=2v_{\mathfrak p}(y)$ is even, so each is even. The values $\delta_i(T_i)$ only
involve primes of $S$.

Finiteness. Every $\alpha$ with even valuations outside $S$ has $(\alpha)=\mathfrak a^2\prod_{\mathfrak p\in S}\mathfrak p^{n_{\mathfrak p}}$. The map
$\alpha\mapsto(n_{\mathfrak p}\bmod2)_{\mathfrak p\in S}$ is well defined on $K(S,2)$, with values in a group of order $2^{|S|}$. On its kernel,
$(\alpha)=\mathfrak b^2$ (put the even powers of the primes of $S$ into $\mathfrak b$), and $\alpha\mapsto[\mathfrak b]$ is a homomorphism to the $2$-torsion of the
class group (changing $\alpha$ by $\gamma^2$ changes $\mathfrak b$ by $(\gamma)$), which is finite (Theorem~\ref{thm:minkbound}). On the kernel of this second
map, $\mathfrak b=(\beta)$, so $\alpha=u\beta^2$ with $u$ a unit: the class of $\alpha$ is that of a unit, and $\mathcal O_K^\times/\mathcal O_K^{\times2}$ is finite
because $\mathcal O_K^\times$ is finitely generated (Theorem~\ref{thm:dirichlet}). So $|K(S,2)|\le2^{|S|}\,|\mathrm{Cl}_K[2]|\,|\mathcal O_K^\times/\mathcal O_K^{\times2}|$.
\end{proof}

\begin{theorem}[weak Mordell--Weil]\label{thm:weakmw}
$E(\Q)/2E(\Q)$ is finite.
\end{theorem}
\begin{proof}
By the two lemmas, $\delta=(\delta_1,\delta_2,\delta_3)$ maps $E(K)/2E(K)$ injectively into the finite group $K(S,2)^3$, so $E(K)/2E(K)$ is finite. It remains to see
that $E(\Q)/2E(\Q)\to E(K)/2E(K)$ has finite kernel $(E(\Q)\cap2E(K))/2E(\Q)$. Let $G=\operatorname{Gal}(K/\Q)$; it acts on $E(K)$ by acting on coordinates, by
group automorphisms (the formulas \eqref{eq:addformula} have rational coefficients). For $P\in E(\Q)\cap2E(K)$ choose $Q_P\in E(K)$ with $2Q_P=P$. For
$\sigma\in G$, $2(\sigma Q_P-Q_P)=\sigma P-P=O$, so $\sigma\mapsto\sigma Q_P-Q_P$ is a map $G\to E[2]$, the four points of order dividing $2$. There are at most
$4^{|G|}$ such maps. If $P$ and $P'$ give the same map, then $Q_P-Q_{P'}$ is fixed by $G$, so it lies in $E(\Q)$, and $P-P'=2(Q_P-Q_{P'})\in2E(\Q)$. So the kernel
has at most $4^{|G|}$ elements.
\end{proof}

\subsection*{Heights}

For a rational number $r=m/n$ in lowest terms, let $H(r)=\max(|m|,|n|)$, and for $P\in E(\Q)$ let $h(P)=\log H(x(P))$, $h(O)=0$. For every $B$, the set
$\{P:h(P)\le B\}$ is finite: finitely many $x$, and two $y$ for each.

\begin{lemma}\label{lem:heightadd}
For every $P_0\in E(\Q)$ there is $c_1$ with $h(P+P_0)\le2h(P)+c_1$ for all $P$.
\end{lemma}
\begin{proof}
We may exclude the finitely many $P$ with $P\in\{O,P_0,-P_0\}$ and assume $P_0\ne O$. Write $x=x(P)=m/e^2$, $y=n/e^3$ as above, so $H(x)=\max(|m|,e^2)$ (the
fraction $m/e^2$ is in lowest terms). From $n^2=m^3+ame^4+be^6$, $|n|\le C\max(|m|,e^2)^{3/2}$ with $C$ depending on $a,b$. The addition formula gives, for $x\ne x_0$,
\[
x(P+P_0)=\Bigl(\frac{y-y_0}{x-x_0}\Bigr)^2-x-x_0=\frac{(xx_0+a)(x+x_0)+2b-2yy_0}{(x-x_0)^2}
\]
(use $y^2=f(x)$, $y_0^2=f(x_0)$ and expand). Multiply numerator and denominator by $e^4$ and by a suitable integer $N$ depending only on $P_0$: they become integers,
$N\bigl[(mx_0+ae^2)(m+x_0e^2)+2be^4-2ney_0\bigr]$ and $N(m-x_0e^2)^2$, of absolute value at most $C'\max(|m|,e^2)^2$, using $|ne|\le C\max(|m|,e^2)^{3/2}e\le
C\max(|m|,e^2)^2$. Cancelling common factors only decreases them, so $H(x(P+P_0))\le C'H(x)^2$.
\end{proof}

\begin{lemma}\label{lem:heightdouble}
There is $c_2$ with $h(2P)\ge4h(P)-c_2$ for all $P\in E(\Q)$.
\end{lemma}
\begin{proof}
The finitely many $P$ with $2P=O$ only affect $c_2$. For $P\notin E[2]$, $x(2P)=\varphi(x)/\psi(x)$ with $\varphi=x^4-2ax^2-8bx+a^2$ and $\psi=4f(x)$, by \eqref{eq:addformula}: $\lambda^2-2x=\frac{f'(x)^2-8xf(x)}{4f(x)}$.
The polynomials $\varphi,\psi$ are coprime, since a common root of $\psi$ is a root $e$ of $f$, where $\varphi(e)=\frac{f'(e)^2}4\ne0$. Homogenise: $\Phi(m,n)=n^4\varphi(m/n)$,
$\Psi(m,n)=n^4\psi(m/n)$, binary forms of degree $4$ with integer coefficients and no common zero in $\mathbb P^1(\C)$ ($\Phi(1,0)=1$). By Lemma~\ref{lem:finite}
(over $\Q$), the ideal $(\Phi,\Psi)$ of $\Q[m,n]$ contains a non-zero polynomial in $m$ alone, and, being generated by forms, each homogeneous part of it: so it
contains a power $m^N$, and likewise $n^N$. Clearing denominators, $R\,m^N$ and $R\,n^N$ are combinations of $\Phi,\Psi$ with integer coefficients, for some
integer $R\ne0$. So for coprime integers $m,n$, $\gcd(\Phi(m,n),\Psi(m,n))$ divides $R$. Next, $(u,v)\mapsto\max(|\Phi(u,v)|,|\Psi(u,v)|)$ is continuous and positive on
the compact set $\max(|u|,|v|)=1$, so it is $\ge c>0$ there, and by homogeneity $\max(|\Phi(m,n)|,|\Psi(m,n)|)\ge c\max(|m|,|n|)^4$. With $x=m/n$ in lowest
terms, $H(x(2P))\ge\max(|\Phi|,|\Psi|)/|R|\ge(c/|R|)H(x)^4$.
\end{proof}

\begin{proof}[Proof of Mordell's theorem]
Let $Q_1,\dots,Q_s$ represent the finitely many classes of $E(\Q)/2E(\Q)$ (Theorem~\ref{thm:weakmw}), and let $c_1$ be the largest constant of
Lemma~\ref{lem:heightadd} for $P_0=-Q_1,\dots,-Q_s$. Given $P$, write $P=Q_{i_1}+2P_1$, $P_1=Q_{i_2}+2P_2$, and so on. Then
\[
\begin{gathered}
4h(P_j)-c_2\le h(2P_j)=h(P_{j-1}-Q_{i_j})\le2h(P_{j-1})+c_1,\\
\text{so}\qquad h(P_j)\le\tfrac12h(P_{j-1})+c,\quad c=\tfrac{c_1+c_2}4.
\end{gathered}
\]
By induction $h(P_j)\le2^{-j}h(P)+2c$, so some $P_j$ lies in the finite set $\mathcal F=\{Q:h(Q)\le2c+1\}$. Unwinding, $P$ is a combination of the $Q_i$ and of an
element of $\mathcal F$: the finite set $\{Q_1,\dots,Q_s\}\cup\mathcal F$ generates $E(\Q)$. The structure theorem for finitely generated abelian groups
(Volume~2) gives the decomposition.
\end{proof}

\begin{remark}
Weil extended the theorem in 1928 to abelian varieties over number fields (\emph{Mordell--Weil}). The height has a better version, the \emph{canonical height}
$\hat h(P)=\lim4^{-n}h(2^nP)$, a positive definite quadratic form on $E(\Q)\otimes\R$, which a direct computation gives. The proof gives no algorithm for $r$: the
descent bounds $|E(\Q)/2E(\Q)|$ from above, and the gap between this bound and the truth is measured by the mysterious Tate--Shafarevich group.
\end{remark}

\section{Congruent numbers}

\begin{proposition}\label{prop:congruent}
Let $n\ge1$ be an integer. The right triangles with rational sides $a,b,c$ and area $n$ (with $a,b>0$) and the points of $E_n:y^2=x^3-n^2x$ are related as follows: the maps
\[
(a,b,c)\mapsto\Bigl(\frac{n(a+c)}b,\ \frac{2n^2(a+c)}{b^2}\Bigr),\qquad(x,y)\mapsto\Bigl(\frac{x^2-n^2}{y},\ \frac{2nx}y,\ \frac{x^2+n^2}{y}\Bigr).
\]
are inverse bijections between the rational triples with $a^2+b^2=c^2$ and $ab=2n$, and the points of $E_n(\Q)$ with $y\ne0$. In particular $n$ is congruent iff
$E_n(\Q)$ has a point with $y\ne0$.
\end{proposition}
\begin{proof}
If $(x,y)\in E_n(\Q)$ with $y\ne0$, the three numbers $a=\frac{x^2-n^2}y$, $b=\frac{2nx}y$, $c=\frac{x^2+n^2}y$ satisfy
$a^2+b^2=\frac{(x^2-n^2)^2+4n^2x^2}{y^2}=c^2$ and $\frac{ab}2=\frac{n(x^3-n^2x)}{y^2}=n$; replacing $a,b$ by $|a|,|b|$ gives a triangle. Conversely, if $a^2+b^2=c^2$ and $ab=2n$, then $b\ne0$, and $a+c\ne0$ (otherwise $b=0$). Put $x=\frac{n(a+c)}b$ and $y=\frac{2n^2(a+c)}{b^2}$. Using
$c^2-b^2=a^2$ and then $2n=ab$ twice,
\[
\begin{aligned}
x^2-n^2&=\frac{n^2\bigl((a+c)^2-b^2\bigr)}{b^2}=\frac{2n^2a(a+c)}{b^2},\\
x^3-n^2x&=\frac{2n^3a(a+c)^2}{b^3}=\frac{n^2a^2(a+c)^2}{b^2}=\frac{4n^4(a+c)^2}{b^4}=y^2.
\end{aligned}
\]
The maps are inverse to each other. Starting from $(x,y)$: $a+c=\frac{2x^2}y$, so $\frac{n(a+c)}b=\frac{2nx^2/y}{2nx/y}=x$, and similarly $y$ comes back. Starting from
$(a,b,c)$: $\frac{x^2-n^2}y=a$ by the first line above, $\frac{2nx}y=b$, and
\[
\frac{x^2+n^2}y=a+\frac{2n^2}y=a+\frac{b^2}{a+c}=\frac{a^2+ac+b^2}{a+c}=c.
\]
\end{proof}

For $n=6$, the $3$-$4$-$5$ triangle gives $P=(12,36)$. Then $2P=\bigl(\frac{25}4,-\frac{35}8\bigr)$ gives the triangle $\frac7{10},\frac{120}7,\frac{1201}{70}$, and
$4P=\bigl(\frac{1442401}{19600},\dots\bigr)$ gives one with hypotenuse $\frac{2094350404801}{241717895860}$ (Figure~\ref{fig:area6}). Each doubling roughly quadruples
the number of digits (Lemma~\ref{lem:heightdouble}), and since $E_6(\Q)$ will turn out to have points of infinite order, there are infinitely many triangles of area
$6$.

\begin{figure}[ht]
\centering\fig{ch12-area6}
\caption{Three right triangles of area $6$, from the points $P$, $2P$, $4P$ on $y^2=x^3-36x$, drawn to the same scale. The legs of the third are
about $1.40$ and $8.55$.}
\label{fig:area6}
\end{figure}

\begin{proposition}\label{prop:entors}
The torsion subgroup of $E_n(\Q)$ is $E_n[2]$, made of $O$, $(0,0)$, $(n,0)$ and $(-n,0)$. Hence $n$ is congruent iff $E_n(\Q)$ has positive rank.
\end{proposition}
\begin{proof}
For a prime $p\equiv3\pmod4$ with $p\nmid2n$, $\tilde E_n(\mathbb F_p)$ has $p+1$ points: $g(x)=x^3-n^2x$ is odd, $-1$ is not a square modulo $p$, so for each of the
$\frac{p-3}2$ pairs $\{x,-x\}$ with $g(x)\ne0$, exactly one of $g(x),g(-x)=-g(x)$ is a non-zero square and gives two points; with the three points of order $2$ and $O$,
the total is $(p-3)+3+1=p+1$. The discriminant is $4(-n^2)^3=-4n^6$, so Proposition~\ref{prop:reduction} applies, and $|T|$ divides $p+1$ for all such $p$, where
$T=E_n(\Q)_{\mathrm{tors}}\supseteq E_n[2]$, of order $4$. If an odd prime $\ell$ divides $|T|$, choose by Dirichlet's theorem a prime $p\nmid2n$ with $p\equiv3\pmod4$
and $p\equiv1\pmod\ell$ (Theorem~\ref{thm:dirichletF} for $K=\Q$): then $\ell\mid p+1\equiv2\pmod\ell$, absurd. If $8\mid|T|$, take $p\equiv3\pmod8$: then $p+1\equiv4\pmod8$.
So $|T|=4$. The last statement follows from Proposition~\ref{prop:congruent}, since the points of order $2$ are those with $y=0$.
\end{proof}

\begin{theorem}[Fermat]\label{thm:onecongruent}
$E_1:y^2=x^3-x$ has rank $0$, with
\[
E_1(\Q)=\{O,(0,0),(1,0),(-1,0)\}.
\]
Hence $1$ is not a congruent number, no right triangle with rational sides has square area, and
the equations $x^4-y^4=z^2$ and $x^4+y^4=z^4$ have no solution in positive integers.
\end{theorem}
\begin{proof}
Here $e_1=0$, $e_2=1$, $e_3=-1$, $K=\Q$, and $S=\{2\}$ (with the sign, $K(S,2)=\{\pm1,\pm2\}$ modulo squares). Since $\delta_1\delta_2\delta_3$ is a square
($=y^2$ for affine points not of order $2$; check the $T_i$), $\delta_3$ is determined by $\delta_1,\delta_2$, and Lemma~\ref{lem:deltahom} shows that
$(\delta_1,\delta_2):E_1(\Q)/2E_1(\Q)\to\{\pm1,\pm2\}^2$ is injective. On the torsion:
\[
\delta(O)=(1,1),\quad\delta(0,0)=(-1,-1),\quad\delta(1,0)=(1,2),\quad\delta(-1,0)=(-1,-2),
\]
using $\delta_1(T_1)=(0-1)(0+1)=-1$, $\delta_2(T_2)=(1-0)(1+1)=2$, and so on. These four values form a subgroup $V$. We show that the image of $\delta$ is $V$.

Let $P=(x,y)$ with $y\ne0$, and $x=d_1u^2$, $x-1=d_2v^2$, $x+1=d_1d_2w^2$ with $d_1,d_2\in\{\pm1,\pm2\}$ and $u,v,w\in\Q^\times$. \emph{Signs.} $x(x-1)(x+1)=y^2>0$, so either
$x>1$ (then $d_1,d_2>0$) or $-1<x<0$ (then $d_1,d_2<0$). \emph{Powers of $2$.} Suppose $|d_1|=2$; then $v_2(x)$ is odd. If $v_2(x)>0$, then $x-1$ and $x+1$ are
$2$-adic units, so $d_2$ (of even valuation) is $\pm1$ and $v_2(x+1)=v_2(d_1d_2)+2v_2(w)$ is odd: absurd. If $v_2(x)<0$, then
$v_2(x-1)=v_2(x+1)=v_2(x)$ is odd, so $|d_2|=2$ and $v_2(x+1)=v_2(d_1d_2)+2v_2(w)$ is even: absurd. So $d_1=\pm1$, with the sign of $d_2$, and $(d_1,d_2)$ is
one of $(1,1),(1,2),(-1,-1),(-1,-2)$: it is in $V$.

So $|E_1(\Q)/2E_1(\Q)|=|V|=4$. But $E_1(\Q)\cong T\times\Z^r$ with $T=(\Z/2)^2$ (Proposition~\ref{prop:entors}), so $E_1(\Q)/2E_1(\Q)\cong(\Z/2)^{2+r}$: $r=0$.

Consequences. By Propositions~\ref{prop:congruent} and~\ref{prop:entors}, $1$ is not congruent. A rational right triangle with area $s^2$ would give, after scaling by $1/s$,
one with area $1$. If $x^4-y^4=z^2$ with $x,y,z>0$, the point $(X,Y)=\bigl(\frac{x^2}{y^2},\frac{xz}{y^3}\bigr)$ satisfies $Y^2=\frac{x^2(x^4-y^4)}{y^6}=X^3-X$ and $Y\ne0$:
impossible. And $x^4+y^4=z^4$ gives $z^4-y^4=(x^2)^2$.
\end{proof}

\begin{beauty}
Fermat's descent for $x^4+y^4=z^2$ (Volume~1) and the proof above are the same argument in different clothes: the map $\delta$ is Fermat's factorisation, and
``smaller'' is the height. What the elliptic curve adds is structure. A failed descent is no longer a dead end but a group $E(\Q)/2E(\Q)$ to be computed,
and the same proof works for every $n$, for every curve, and, with Weil's generalisation, for curves of any genus inside their Jacobians.
\end{beauty}

\begin{remark}[Birch--Swinnerton-Dyer, Tunnell]
Running the descent for many curves, Birch and Swinnerton-Dyer noticed in the early 1960s that the rank seemed to be governed by the numbers $N_p=|\tilde E(\mathbb F_p)|$:
$\prod_{p\le X}N_p/p$ grows like $(\log X)^r$. In terms of the $L$-function $L(E,s)=\prod_p(1-a_pp^{-s}+p^{1-2s})^{-1}$, $a_p=p+1-N_p$ (for good $p$), their conjecture
says that $L(E,s)$ has a zero of order exactly $r$ at $s=1$. It is one of the Clay Millennium problems. For $E_n$, Tunnell (1983) computed $L(E_n,1)$ with modular forms
of weight $\frac32$: if $n$ is odd and squarefree and congruent, then the number of solutions of $n=2x^2+y^2+32z^2$ is half the number of solutions of $n=2x^2+y^2+8z^2$; and
the converse holds if the conjecture holds for $E_n$. Every $n\equiv5,6,7\pmod8$ should be congruent; this is known for primes $p\equiv5,7\pmod8$ (Heegner, Monsky).
\end{remark}

\section{Exercises}

\begin{exercise}[$\star$]\label{ex:x3plus1}
Show that the torsion subgroup of $y^2=x^3+1$ is cyclic of order $6$, generated by $(2,3)$, using Theorem~\ref{thm:nagelllutz} alone.
\end{exercise}

\begin{exercise}[$\star$]\label{ex:five}
Show that $5$ and $7$ are congruent: find points on $E_5$ and $E_7$ (try $x=-4$ and $x=25$) and the corresponding triangles.
\end{exercise}

\begin{exercise}[$\star\star$]\label{ex:flex}
Show that a point $P\ne O$ of $E$ has order $3$ iff it is a flex, that is, the tangent at $P$ meets $E$ at $P$ with multiplicity $3$. Find the points of order $3$ of
$y^2=x^3+1$.
\end{exercise}

\begin{exercise}[$\star\star$]\label{ex:bachet}
(Bachet, 1621.) On $y^2=x^3+c$, show that $x(2P)=\dfrac{x^4-8cx}{4y^2}$. Starting from $3^3=27=5^2+2$, find a second rational solution of $y^2=x^3-2$, and
show that the point $(3,5)$ has infinite order.
\end{exercise}

\begin{exercise}[$\star\star$]\label{ex:supersing}
Let $p\equiv2\pmod3$ be a prime and $c\not\equiv0\pmod p$. Show that $y^2=x^3+c$ has exactly $p+1$ points over $\mathbb F_p$ (including $O$). Deduce that the torsion
subgroup of $y^2=x^3+c$ over $\Q$ has order dividing $6$ when $c$ is an integer.
\end{exercise}

\begin{exercise}[$\star\star$]\label{ex:torsion7}
Show that $P=(3,8)$ on $y^2=x^3-43x+166$ has order $7$ and that it generates the torsion subgroup. (Compute $2P,3P,4P$, and count points modulo $3$.)
\end{exercise}

\begin{exercise}[$\star\star\star$]\label{ex:twonotcong}
Show that $2$ is not a congruent number, by adapting the proof of Theorem~\ref{thm:onecongruent} to $y^2=x^3-4x$.
\end{exercise}

\section{Hints}

\begin{hint}{ex:x3plus1}
$D=27$: the possible $y$ are $0,\pm1,\pm3$. Then compute multiples of $(2,3)$.
\end{hint}
\begin{hint}{ex:five}
Use the inverse map of Proposition~\ref{prop:congruent}.
\end{hint}
\begin{hint}{ex:flex}
$3P=O$ iff $2P=-P$ iff $P*P=P$.
\end{hint}
\begin{hint}{ex:bachet}
Use \eqref{eq:addformula} with $a=0$. For infinite order, compute $2P$ and use Theorem~\ref{thm:nagelllutz}.
\end{hint}
\begin{hint}{ex:supersing}
When $p\equiv2\pmod3$, $x\mapsto x^3$ is a bijection of $\mathbb F_p$. For the deduction, use Proposition~\ref{prop:reduction} with two primes.
\end{hint}
\begin{hint}{ex:torsion7}
$3P$ and $4P$ have the same $x$. Modulo $3$ the curve is $y^2=x^3-x+1$.
\end{hint}
\begin{hint}{ex:twonotcong}
Here $e_i=0,2,-2$, and $\delta(T)$ for the three points of order $2$ are $(-1,-2)$, $(2,2)$, $(-2,-1)$. Exclude the other pairs of the right signs by looking at
$v_2(x)$.
\end{hint}

\section{Solutions}

\begin{solution}{ex:x3plus1}
\idea List the Nagell--Lutz candidates, then check orders with the addition formula.
\writeup $D=4\cdot0+27=27$, so a torsion point has $y=0$ or $y^2\mid27$: $y\in\{0,\pm1,\pm3\}$. Then $x^3=y^2-1\in\{-1,0,8\}$: the points $(-1,0)$, $(0,\pm1)$, $(2,\pm3)$.
For $P=(2,3)$: $\lambda=\frac{3\cdot4}{6}=2$, $x(2P)=4-4=0$, $y(2P)=-3-2(0-2)=1$: $2P=(0,1)$. Then $3P=P+2P$: $\lambda=\frac{1-3}{0-2}=1$, $x_3=1-2-0=-1$, $y_3=-3-1\cdot(-1-2)=0$:
$3P=(-1,0)$, of order $2$. So $P$ has order $6$, and its six multiples $O,(2,3),(0,1),(-1,0),(0,-1),(2,-3)$ exhaust the candidates; all are torsion, so the torsion
subgroup is $\langle P\rangle\cong\Z/6$.
\end{solution}

\begin{solution}{ex:five}
\idea Points are easy to check; the triangle comes from the formulas of Proposition~\ref{prop:congruent}.
\writeup On $y^2=x^3-25x$: $x=-4$ gives $-64+100=36$, $y=6$. Then $a=\frac{16-25}6=-\frac32$, $b=\frac{2\cdot5\cdot(-4)}6=-\frac{20}3$, $c=\frac{16+25}6=\frac{41}6$: the
triangle $\frac32,\frac{20}3,\frac{41}6$, Fibonacci's, with $\frac94+\frac{400}9=\frac{1681}{36}$ and area $\frac12\cdot\frac32\cdot\frac{20}3=5$. On $y^2=x^3-49x$: $x=25$ gives
$15625-1225=14400=120^2$. Then $a=\frac{625-49}{120}=\frac{24}5$, $b=\frac{2\cdot7\cdot25}{120}=\frac{35}{12}$, $c=\frac{674}{120}=\frac{337}{60}$, and $\frac12\cdot\frac{24}5\cdot\frac{35}{12}=7$.
\end{solution}

\begin{solution}{ex:flex}
\idea Translate $3P=O$ into a statement about the tangent line.
\writeup $3P=O$ iff $2P=-P$. Now $2P=O*(P*P)$ and $-P=O*P$, and $R\mapsto O*R$ is injective (it is $(x,y)\mapsto(x,-y)$), so $2P=-P$ iff $P*P=P$, that is, iff the
tangent at $P$ meets $E$ a third time at $P$: the intersection multiplicity at $P$ is $3$. (If $2P=O$, the tangent is vertical and $P*P=O\ne P$.) For $y^2=x^3+1$:
at $(0,\pm1)$ the tangent is horizontal ($\lambda=\frac{3\cdot0}{\pm2}=0$), the line $y=\pm1$, and $x^3+1-1=x^3$ has a triple root at $0$: these are flexes, of order $3$.
Any other point of order $3$ would be torsion, hence one of the six points of Exercise~\ref{ex:x3plus1}; only $(0,\pm1)$ have order $3$.
\end{solution}

\begin{solution}{ex:bachet}
\idea Doubling is the tangent construction; Nagell--Lutz detects infinite order through a non-integral multiple.
\writeup With $a=0$, $\lambda=\frac{3x^2}{2y}$ and $x(2P)=\frac{9x^4}{4y^2}-2x=\frac{9x^4-8xy^2}{4y^2}=\frac{9x^4-8x(x^3+c)}{4y^2}=\frac{x^4-8cx}{4y^2}$. For $c=-2$ and $(3,5)$:
$x(2P)=\frac{81+48}{100}=\frac{129}{100}$, and $y(2P)=-5-\frac{27}{10}\bigl(\frac{129}{100}-3\bigr)=-5+\frac{27\cdot171}{1000}=\frac{-5000+4617}{1000}=-\frac{383}{1000}$. Check:
$\bigl(\frac{383}{1000}\bigr)^2=\frac{146689}{10^6}$ and $\bigl(\frac{129}{100}\bigr)^3-2=\frac{2146689-2000000}{10^6}$. Since $2P$ has non-integral coordinates, it is not torsion
(Theorem~\ref{thm:nagelllutz}), so neither is $P$.
\end{solution}

\begin{solution}{ex:supersing}
\idea A bijective cube map makes $x^3+c$ run through $\mathbb F_p$ exactly once.
\writeup As $p\equiv2\pmod3$, $\gcd(3,p-1)=1$, so $x\mapsto x^3$ is a bijection of $\mathbb F_p$ (its kernel on $\mathbb F_p^\times$ is trivial), and so is $x\mapsto x^3+c$.
The number of affine points is $\sum_{u\in\mathbb F_p}\#\{y:y^2=u\}=\sum_u\bigl(1+\bigl(\frac up\bigr)\bigr)=p$, as the Legendre symbol sums to $0$. With $O$: $p+1$. For
$c\in\Z$, $c\ne0$, $D=27c^2$, so by Proposition~\ref{prop:reduction} $|T|$ divides $p+1$ for every prime $p\equiv2\pmod3$ with $p\nmid6c$. If a prime $\ell\ge5$ divided $|T|$, Dirichlet's theorem gives such a
$p$ with $p\equiv1\pmod\ell$ (a condition modulo $3\ell$), and then $\ell\mid p+1\equiv2$: absurd. If $4\mid|T|$, take $p\equiv5\pmod{12}$, so $p+1\equiv6\pmod{12}$ is not divisible by $4$. If $9\mid|T|$, take $p\equiv2\pmod9$:
$p+1\equiv3\pmod9$. So $|T|$ divides $6$.
\end{solution}

\begin{solution}{ex:torsion7}
\idea Find the order by computing multiples; bound the torsion with a reduction.
\writeup With $a=-43$: at $P=(3,8)$, $\lambda=\frac{27-43}{16}=-1$, $x(2P)=1-6=-5$, $y(2P)=-8+1\cdot(-5-3)=-16$: $2P=(-5,-16)$. Then $P+2P$: $\lambda=\frac{-16-8}{-5-3}=3$,
$x=9-3+5=11$, $y=-8-3(11-3)=-32$: $3P=(11,-32)$. And $4P=2(2P)$: $\lambda=\frac{75-43}{-32}=-1$, $x=1+10=11$, $y=16+1\cdot(11+5)=32$: $4P=(11,32)=-3P$. So $7P=O$ and $P$ has
order $7$. Here $D=4(-43)^3+27\cdot166^2=425984=2^{15}\cdot13$, so $p=3$ is a good prime, and $\tilde E:y^2=x^3-x+1$ over $\mathbb F_3$: $x^3-x+1=1$ for all $x\in\mathbb F_3$,
so each $x$ gives two points: $6+1=7$ points. Hence $|T|$ divides $7$ (Proposition~\ref{prop:reduction}), and $T=\langle P\rangle$.
\end{solution}

\begin{solution}{ex:twonotcong}
\idea Same descent. The image of $\delta$ lies between the torsion image $V$ and the group $\Sigma$ of pairs with equal signs, and $[\Sigma:V]=2$; so it is enough to
exclude a single pair, and a congruence modulo $8$ does it.
\writeup For $y^2=x^3-4x$: $e_1=0$, $e_2=2$, $e_3=-2$, $S=\{2\}$, and $(\delta_1,\delta_2)$ is injective on $E_2(\Q)/2E_2(\Q)$ with values in $\{\pm1,\pm2\}^2$. On the torsion:
$\delta(0,0)=\bigl((0-2)(0+2),\,0-2\bigr)=(-4,-2)\equiv(-1,-2)$; $\delta(2,0)=(2,(2-0)(2+2))=(2,8)\equiv(2,2)$; $\delta(-2,0)=(-2,-4)\equiv(-2,-1)$. With $(1,1)$ these form
a subgroup $V$ of order $4$. For $P=(x,y)$ with $y\ne0$, $x(x-2)(x+2)>0$ forces $x>2$ or $-2<x<0$, so $\delta_1(P)$ and $\delta_2(P)=x-2$ have the same sign. Hence the
image of $\delta$ lies in the subgroup $\Sigma$ of pairs with equal signs, of order $8$, and contains $V$. It is $V$ unless it contains $(1,2)\in\Sigma\setminus V$.

Suppose $\delta(P)=(1,2)$; then $y\ne0$ (points of order $2$ map into $V$), and $x=u^2$, $x-2=2v^2$, $x+2=2w^2$ with $u,v,w\in\Q$. Write $v=V/N$, $w=W/N$ with
$N\ge1$ minimal, so $\gcd(V,W,N)=1$. From $w^2-v^2=2$, $W^2-V^2=2N^2$; from $u^2=2v^2+2$, $(uN)^2=2(V^2+N^2)$, so $uN$ is an integer. If $N$ is odd, $W^2-V^2\equiv2\pmod8$,
impossible since squares are $0,1,4$ modulo $8$ and no two differ by $2$. If $N$ is even, $W^2\equiv V^2\pmod8$, so $W,V$ have the same parity; both even contradicts
$\gcd(V,W,N)=1$, so both are odd, and then $2(V^2+N^2)\equiv2\pmod8$ is not a square. So $(1,2)$ is not in the image, the image is $V$, and $E_2(\Q)/2E_2(\Q)$ has order $4$.
As in Theorem~\ref{thm:onecongruent} (the torsion is $E_2[2]$ by Proposition~\ref{prop:entors}), the rank is $0$, and $2$ is not congruent.
\end{solution}

\section*{Sources and further reading}

\begin{itemize}
\item J.~H.~Silverman and J.~Tate, \emph{Rational Points on Elliptic Curves}, 2nd ed., Springer UTM, 2015. The filtration, Nagell--Lutz, heights and Mordell's
theorem with rational $2$-torsion; the friendliest introduction.
\item J.~H.~Silverman, \emph{The Arithmetic of Elliptic Curves}, 2nd ed., Springer GTM~106, 2009, Chapters~III, VII, VIII and~X. Weak Mordell--Weil over number fields,
the Kummer map, descent in general.
\item N.~Koblitz, \emph{Introduction to Elliptic Curves and Modular Forms}, 2nd ed., Springer GTM~97, 1993. The congruent number problem from Fibonacci to Tunnell.
\item L.~C.~Washington, \emph{Elliptic Curves: Number Theory and Cryptography}, 2nd ed., CRC, 2008. The explicit halving argument of Lemma~\ref{lem:deltahom}, and
cryptography.
\item J.~B.~Tunnell, ``A classical Diophantine problem and modular forms of weight $3/2$'', \emph{Invent.\ Math.} 72 (1983), 323--334.
\item B.~Mazur, ``Modular curves and the Eisenstein ideal'', \emph{Publ.\ Math.\ IH\'ES} 47 (1977), 33--186.
\item B.~J.~Birch and H.~P.~F.~Swinnerton-Dyer, ``Notes on elliptic curves II'', \emph{J.\ reine angew.\ Math.} 218 (1965), 79--108.
\item \emph{History.} The Arab manuscript of the tenth century, Fibonacci's \emph{Liber quadratorum} (1225) and Fermat's
descent: L.~E.~Dickson, \emph{History of the Theory of Numbers}, vol.~II, Carnegie Institution, 1920, Chapter~XVI; and
Koblitz's book above. Poincar\'e's question: H.~Poincar\'e, ``Sur les propri\'et\'es arithm\'etiques des courbes
alg\'ebriques'', \emph{J.\ Math.\ Pures Appl.}\ (5) 7 (1901), 161--233.
\end{itemize}
